%! Rates of Change in Polar Functions — Unit 7, Lesson 7.8 — Slides (Beamer)
\documentclass[aspectratio=169, 11pt]{beamer}
\usepackage{precalculus-beamer}   % provides \navyheader{} and \sectionlabel[color]{}

\begin{document}

% TITLE SLIDE
{
\setbeamercolor{background canvas}{bg=navy}
\begin{frame}[plain]
  \vspace{2.8cm}\hspace{0.55cm}%
  \begin{minipage}{0.9\textwidth}
    {\color{goldacc}\bfseries\small LESSON 3.15}\\[0.5cm]
    {\color{white}\bfseries\LARGE Rates of Change in Polar Functions}\\[1.2cm]
    {\color{skymid}\small Precalculus~$\cdot$~Unit 3}
  \end{minipage}
\end{frame}
}

%----------------------------------------------------------------------
% Frame 1 — Average Rate of Change of r
%----------------------------------------------------------------------
\begin{frame}[t, plain]
  \navyheader{Average Rate of Change of $r = f(\theta)$}

  \begin{block}{Formula (EK 3.15.A.1)}
    $$\text{ARC} = \frac{f(\theta_2) - f(\theta_1)}{\theta_2 - \theta_1}$$
    Measures how quickly $r$ (the radius) changes per \textbf{radian} of angle change.
  \end{block}

  \vspace{0.3cm}
  \textbf{Example:} $r = 2\sin\theta$

  \vspace{0.2cm}
  \begin{tabular}{|c|c|c|}
    \hline
    $\theta$ & $0$ & $\pi/2$ \\
    \hline
    $r = 2\sin\theta$ & $0$ & $2$ \\
    \hline
  \end{tabular}

  \vspace{0.3cm}
  ARC from $\theta = 0$ to $\theta = \pi/2$:
  $$\frac{2 - 0}{\pi/2} = \frac{4}{\pi} \approx 1.27 \text{ per radian}$$

  \vspace{0.2cm}
  \begin{itemize}
    \item Positive ARC $\Rightarrow$ $r$ is \textbf{increasing}
    \item Negative ARC $\Rightarrow$ $r$ is \textbf{decreasing}
  \end{itemize}
\end{frame}

%----------------------------------------------------------------------
% Frame 2 — The Four Cases
%----------------------------------------------------------------------
\begin{frame}[t, plain]
  \navyheader{Four Cases: Sign $\times$ Direction (EK 3.15.A.2--3)}

  \begin{block}{Classification Table}
    \renewcommand{\arraystretch}{1.5}
    \begin{tabular}{|l|l|l|}
      \hline
      \textbf{Sign of $r$} & \textbf{$r$ is \ldots} & \textbf{Curve behavior in plane} \\
      \hline
      $r > 0$ & Increasing & Moving \textbf{away} from origin \\
      $r > 0$ & Decreasing & Moving \textbf{toward} origin \\
      $r < 0$ & Increasing (less negative) & Toward origin (opposite ray) \\
      $r < 0$ & Decreasing (more negative) & Away from origin (opposite ray) \\
      \hline
    \end{tabular}
  \end{block}

  \vspace{0.4cm}
  \begin{block}{Key Idea: Negative $r$}
    When $r < 0$, the point is plotted \textbf{opposite} the terminal ray at
    distance $|r|$ from the origin. A ``less negative'' $r$ means the point
    is moving \emph{toward} the origin on that opposite ray.
  \end{block}
\end{frame}

%----------------------------------------------------------------------
% Frame 3 — Reading Intervals from a Table
%----------------------------------------------------------------------
\begin{frame}[t, plain]
  \navyheader{Worked Example: Reading Intervals from a Table}

  \textbf{Given:} $r = 1 + \cos\theta$

  \vspace{0.2cm}
  \renewcommand{\arraystretch}{1.3}
  \begin{tabular}{|c|c|c|c|c|c|c|}
    \hline
    $\theta$ & $0$ & $\pi/4$ & $\pi/2$ & $3\pi/4$ & $\pi$ & $5\pi/4$ \\
    \hline
    $r$ & $2$ & $\approx 1.71$ & $1$ & $\approx 0.29$ & $0$ & $\approx 0.29$ \\
    \hline
  \end{tabular}

  \vspace{0.4cm}
  \textbf{Question:} Classify $[0, \pi/2]$ and $[\pi, 5\pi/4]$.

  \vspace{0.2cm}
  \begin{itemize}
    \item $[0, \pi/2]$: $r$ goes from 2 to 1 --- positive and \textbf{decreasing} (P\&D)
    \item[$\Rightarrow$] Curve moves \textbf{toward} the origin.
    \item[$\,$]
    \item $[\pi, 5\pi/4]$: $r$ goes from 0 to $0.29$ --- positive and \textbf{increasing} (P\&I)
    \item[$\Rightarrow$] Curve moves \textbf{away} from the origin.
  \end{itemize}

  \vspace{0.2cm}
  \begin{block}{Strategy}
    Compare successive table values. Increasing values $\Rightarrow$ P\&I or N\&I;
    decreasing values $\Rightarrow$ P\&D or N\&D. Check sign of $r$ to assign P or N.
  \end{block}
\end{frame}

%----------------------------------------------------------------------
% Frame 4 — Connecting to the Curve in the Plane
%----------------------------------------------------------------------
\begin{frame}[t, plain]
  \navyheader{Connecting Rate of Change to Curve Shape}

  \begin{block}{What does the rate of change tell us?}
    The sign and magnitude of the ARC tell us \emph{how} the polar curve is being
    traced as $\theta$ sweeps from 0 to $2\pi$.
  \end{block}

  \vspace{0.3cm}
  \textbf{Cardioid $r = 1 + \cos\theta$:}
  \begin{itemize}
    \item On $[0, \pi]$: $r$ decreases from 2 to 0 $\Rightarrow$ P\&D $\Rightarrow$ spirals toward the pole.
    \item On $[\pi, 2\pi]$: $r$ increases from 0 to 2 $\Rightarrow$ P\&I $\Rightarrow$ spirals outward.
    \item The curve traces a heart shape, with the point touching the origin at $\theta = \pi$.
  \end{itemize}

  \vspace{0.3cm}
  \textbf{Limaçon with inner loop $r = 1 - 2\sin\theta$:}
  \begin{itemize}
    \item When $r < 0$: the curve is traced on the \emph{opposite} ray --- this produces the inner loop.
    \item N\&D portions move the traced point farther from the origin on the opposite ray.
  \end{itemize}

  \vspace{0.3cm}
  \begin{block}{Big Idea}
    Classifying intervals by sign and direction gives a complete description of
    how any polar curve is traced --- no graphing required.
  \end{block}
\end{frame}

\end{document}
